\section{Sample Importance and Distributional Selection} \label{sec:27-background-importanceselection}

\begin{foldable}
  This section develops the two technical pillars of the TAKE method (Chapter~\ref{ch:50-research03}), parallel in role to \S\ref{sec:25-background-synthetic} for Track~A.
  Where \S\ref{sec:25-background-synthetic} addressed distributional coverage at the \emph{generative} level --- how to ensure a GAN produces images that span the target distribution --- this section addresses coverage at the \emph{selection} level: given a large pool of candidate examples with associated importance scores, how should a compact subset be chosen to minimise distributional divergence from the reweighted source?
\end{foldable}

\subsection{Influence Functions and Hard-Sample Bias}
\label{sec:27-background-importanceselection-influence}

\begin{foldable}
  \textbf{Classical influence functions.}
  Influence functions~\cite{koh2017understanding} formalise the intuition that different training examples contribute differently to a model's predictions.
  The influence of a training example $z_n = (\mathbf{x}_n, y_n)$ on the model's loss at a test point $z_{\text{test}}$ is approximated using implicit differentiation:
  \begin{equation}
    \mathcal{I}(z_n, z_{\text{test}}) \approx
    -\nabla_\theta \mathcal{L}(z_{\text{test}}, \hat\theta)^\top \,
    H_{\hat\theta}^{-1} \,
    \nabla_\theta \mathcal{L}(z_n, \hat\theta),
    \label{eq:influence}
  \end{equation}
  where $\hat\theta$ is the converged model and $H_{\hat\theta} = \frac{1}{N}\sum_n \nabla^2_\theta \mathcal{L}(z_n, \hat\theta)$ is the Hessian of the training loss.
  A large positive influence indicates that removing $z_n$ from the training set would increase the test loss --- i.e., $z_n$ is informative.
  A large negative influence indicates that $z_n$ hurts generalization.
\end{foldable}

\begin{foldable}
  The practical limitation of classical influence functions is computational. Hessian inversion is $O(p^2)$ in the number of parameters $p$, making it intractable for models of even moderate scale.
  For BERT-base with $p \approx 110$M parameters, storing a dense Hessian requires several terabytes of memory.
\end{foldable}

\begin{foldable}
  \textbf{Tractable proxies.}
  TracIn~\cite{pruthi2020estimating} replaces Hessian inversion with a sum of gradient dot-products across training checkpoints:
  \begin{equation}
    \mathrm{TracIn}(z_n, z_{\text{test}}) = \sum_{t=1}^{T} \eta_t \,
    \nabla_\theta \mathcal{L}(z_{\text{test}}, \theta_t)^\top \,
    \nabla_\theta \mathcal{L}(z_n, \theta_t),
    \label{eq:tracin}
  \end{equation}
  where $\eta_t$ is the learning rate and $\theta_t$ is the checkpoint at step $t$.
  TracIn requires no Hessian and is naturally trajectory-aware: it measures how much $z_n$ contributed to the model's improvement throughout training, not just at convergence.
  TRAK~\cite{park2023trak} scales influence estimation further by projecting per-sample gradients onto a random low-dimensional subspace before computing dot-products, reducing memory cost to $O(p \cdot r)$ where $r \ll p$ is the projection dimension.
\end{foldable}

\begin{foldable}
  A simpler proxy --- the gradient norm $\|\nabla_\theta \mathcal{L}(z_n, \hat\theta)\|$ evaluated at the converged checkpoint --- is widely used as a per-sample importance score due to its near-zero computational overhead.
  The gradient norm at convergence measures how far example $z_n$ is from being fully fitted by the current model.
\end{foldable}

\begin{foldable}
  \textbf{Hard-sample bias.}
  The gradient norm proxy, and to a lesser extent any single-checkpoint influence estimator, suffers from a systematic bias at convergence.
  After sufficient training, a well-fitted model assigns near-zero loss to clean, easy examples; their gradient norms are correspondingly small.
  Noisy, mislabelled, or atypical examples, however, resist convergence: the model cannot simultaneously fit them and the rest of the dataset, so their gradient norms remain large throughout training.
  A selector that maximises gradient norms at convergence therefore preferentially selects noisy and hard examples --- precisely the examples that are most harmful to a student trained on the resulting surrogate.
\end{foldable}

\begin{foldable}
  Dataset cartography~\cite{swayamdipta2020dataset} provides empirical evidence for this mechanism: when training dynamics are traced over epochs using per-example confidence and variability, the ``hard-to-learn'' stratum (low confidence, high variability) is strongly enriched in mislabelled examples.
  Co-teaching~\cite{han2018co} arrives at the same conclusion from a noise-robust learning perspective: at late training stages, the large-loss examples that cross-network selection identifies as ``disagreed'' are predominantly noisy.
\end{foldable}

\begin{foldable}
  \textbf{Trajectory integration as correction.}
  If single-checkpoint gradient norms are biased by late-stage noise dominance, averaging gradient norms across $T$ checkpoints from different stages of training dilutes this bias.
  Early-stage checkpoints produce large gradients for \emph{all} examples, not just noisy ones; late-stage checkpoints produce large gradients mainly for noisy examples.
  Their average assigns moderate, calibrated scores to the full spectrum of the training set.
  Chapter~\ref{ch:50-research03} formalises this trajectory-integrated score as $\kappa_n = \frac{1}{T}\sum_{t=1}^{T} \|\nabla_\theta \mathcal{L}(z_n, \theta_t)\|$ and proves in Proposition~2 that $\kappa_n$ is an unbiased estimator of example importance under mild conditions on the training trajectory.
\end{foldable}

\begin{foldable}
  Correcting the importance score is necessary but not sufficient for distributional coverage.
  The selection rule applied to the corrected scores determines whether the chosen subset covers the reweighted source distribution or concentrates on a high-scoring subset of it.
\end{foldable}

\subsection{Optimal Transport and Distributional Coverage}
\label{sec:27-background-importanceselection-ot}

\begin{foldable}
  \textbf{Wasserstein distance.}
  The Wasserstein-$p$ distance~\cite{villani2008optimal} between two probability distributions $\mu$ and $\nu$ on a metric space $(\mathcal{Z}, d)$ is defined as:
  \begin{equation}
    \label{eq:wasserstein}
    W_p(\mu, \nu) = \left(\inf_{\gamma \in \Gamma(\mu,\nu)} \int d(z, z')^p \, \mathrm{d}\gamma(z, z') \right)^{1/p},
  \end{equation}
  where $\Gamma(\mu, \nu)$ is the set of all couplings (joint distributions) with marginals $\mu$ and $\nu$.
  Geometrically, $W_p(\mu, \nu)$ is the minimum cost of transporting mass from $\mu$ to $\nu$, where cost is measured by distance in the ambient space.
  Unlike KL and Jensen-Shannon divergences, the Wasserstein distance is well-defined even when $\mu$ and $\nu$ have disjoint supports --- a critical property when $\mu$ is the full training distribution and $\nu$ is a sparse surrogate.
\end{foldable}

\begin{foldable}
  \textbf{Discrete optimal transport.}
  In the dataset distillation setting, both the source (full training set with importance weights) and the target (synthetic candidate pool) are discrete.
  Let the source be $\{(z_n, w_n)\}_{n=1}^{N}$ where $w_n = \kappa_n / \sum_{n'} \kappa_{n'}$ are the normalised trajectory-integrated importance weights.
  Let the candidate pool be $\{(\tilde{z}_c, 1/C)\}_{c=1}^{C}$ with uniform weights.
  The discrete OT problem seeks a transport plan $\gamma^* \in \mathbb{R}^{N \times C}_{\geq 0}$ that minimises the total transport cost:
  \begin{equation}
    \gamma^* = \arg\min_{\gamma \geq 0} \langle C_{\mathrm{dist}}, \gamma \rangle
    \quad \text{s.t.} \quad \gamma \mathbf{1}_C = \mathbf{w}, \quad
    \gamma^\top \mathbf{1}_N = \tfrac{1}{C}\mathbf{1}_C,
    \label{eq:discrete-ot}
  \end{equation}
  where $C_{\mathrm{dist}} \in \mathbb{R}^{N \times C}$ is the pairwise cost matrix with $C_{\mathrm{dist},nc} = d(z_n, \tilde{z}_c)^2$.
  The entry $\gamma^*_{nc}$ encodes how much importance mass from source example $z_n$ is assigned to candidate $\tilde{z}_c$.
  The subset of $K$ candidates with the largest total assigned mass $\sum_n \gamma^*_{nc}$ is selected as the surrogate $\tilde{\mathcal{D}}$.
\end{foldable}

\begin{foldable}
  \textbf{Sinkhorn algorithm.}
  The discrete OT problem in Equation~\ref{eq:discrete-ot} is a linear programme with $N \times C$ variables, which is expensive to solve exactly for large $N$ and $C$.
  The Sinkhorn algorithm~\cite{cuturi2013sinkhorn} adds an entropic regularization term $\varepsilon \, \mathrm{KL}(\gamma \| \mathbf{w} \otimes \tfrac{1}{C}\mathbf{1}_C)$ to the objective, converting the LP to a strongly convex problem with a closed-form solution via alternating row and column normalization.
  The resulting regularised plan $\gamma_\varepsilon^*$ converges to $\gamma^*$ as $\varepsilon \to 0$ and can be computed in $O(N \cdot C)$ per iteration, making it practical for large candidate pools.
  The regularization parameter $\varepsilon$ controls the smoothness of the transport plan: large $\varepsilon$ produces diffuse plans that spread mass broadly, small $\varepsilon$ produces concentrated plans closer to the exact OT solution.
\end{foldable}

\begin{foldable}
  \textbf{Why OT outperforms greedy and clustering-based selection.}
  $k$-means clustering (used in DaLLME~\cite{tao2024textual}) minimises within-cluster variance without regard to the source importance distribution: clusters in high-density but low-importance regions are treated the same as clusters in high-importance regions.
  Greedy $k$-centre selection minimises the maximum distance from any source point to the nearest selected candidate, again ignoring importance weighting.
  Optimal transport, by contrast, minimises the expected transport cost from the \emph{importance-reweighted} source distribution to the selected surrogate.
  When the source weights $w_n = \kappa_n / \sum_{n'}\kappa_{n'}$ reflect the trajectory-integrated importance scores, OT selection directly minimises the distributional divergence between the surrogate and the semantically important regions of the full training distribution.
\end{foldable}

\begin{foldable}
  \textbf{Toy illustration.}
  A concrete small example makes the contrast between the three selection strategies precise.
  Suppose the full training set contains six examples partitioned into two semantic regions: Region~A holds four examples, each assigned a trajectory-integrated importance weight of~$0.25$ after normalization, giving Region~A a combined weight of~$1.0$;
  Region~B holds two examples, each with weight~$0.05$, giving Region~B a combined weight of~$0.10$.
  Region~A therefore carries roughly ten times the importance mass of Region~B.
  From a candidate pool of four synthetic examples --- two near Region~A and two near Region~B --- the task is to select $K = 2$ candidates as the surrogate.
\end{foldable}

\begin{foldable}
  $k$-means partitions the four candidates into two clusters based on geometric proximity in feature space, without consulting the importance weights.
  If the two Region-B candidates are geometrically cohesive (e.g., both lie in a low-frequency corner of the embedding space), they form one cluster; the two Region-A candidates form the other.
  $k$-means then selects one representative centroid per cluster, yielding one Region-A candidate and one Region-B candidate.
  The resulting surrogate allocates exactly half its capacity to Region~B despite Region~B carrying only~$9\%$ of the total importance mass --- a factor-of-five mismatch relative to the importance-weighted distribution.
\end{foldable}

\begin{foldable}
  Greedy $k$-centre selection chooses the candidate set that minimises the maximum distance from any source point to its nearest selected candidate, again without reference to importance weights.
  In the toy setting, the two geometric extremes of the candidate pool are one Region-A candidate and one Region-B candidate; selecting them minimises the worst-case coverage radius.
  The outcome matches $k$-means: one representative per region, irrespective of the ten-to-one importance imbalance.
  Both geometry-based methods conflate geometric coverage with distributional fidelity, and when the two are decoupled --- as they routinely are in text datasets where low-importance examples cluster in distinctive surface-form niches --- geometric coverage actively misleads the selection.
\end{foldable}

\begin{foldable}
  \ac{OT} with the trajectory-integrated weights $\mathbf{w} = (0.25, 0.25, 0.25, 0.25, 0.05, 0.05)$ constructs a transport plan that must move all of Region~A's mass ($1.0$ in unnormalised units) and all of Region~B's mass ($0.10$) to the candidate pool.
  Because the two Region-A candidates are the nearest feasible receivers for Region~A's mass, the Sinkhorn plan assigns approximately~$0.50$ units of mass to each Region-A candidate and approximately~$0.05$ units to each Region-B candidate.
  Selecting the $K = 2$ candidates with the largest total assigned mass therefore returns both
  Region-A candidates and neither Region-B candidate.
  The surrogate covers Region~A fully and omits Region~B entirely --- which is precisely the outcome that minimises the Wasserstein distance from the importance-reweighted source distribution to the selected subset.
\end{foldable}

\begin{foldable}
  The punchline is that \ac{OT}'s advantage over greedy and clustering-based methods is not about geometric coverage: the \ac{OT} surrogate in this example actually has worse worst-case geometric coverage than either $k$-means or $k$-centre, because it leaves Region~B unrepresented.
  The advantage is about \emph{importance-weighted} coverage.
  When the source weights reflect training dynamics --- as the trajectory-integrated scores $\kappa_n$ are designed to do --- \ac{OT} selection concentrates surrogate capacity on the regions that matter most to the learning objective, and the distributional divergence between the surrogate and the reweighted source is minimised in a principled, quantifiable sense.
\end{foldable}

\begin{foldable}
  \textbf{Composition in Chapter~\ref{ch:50-research03}.}
  The TAKE pipeline composes the two subsections of this section directly.
  The trajectory-integrated scores $\kappa_n$ (\S\ref{sec:27-background-importanceselection-influence}) provide the importance-reweighted marginal $\mathbf{w}$ for the discrete OT problem.
  The Sinkhorn algorithm then selects the $K$ synthetic candidates whose combined assigned mass is largest, formally minimising the distributional gap between the surrogate and the importance-weighted training distribution.
  Theorem~1 in \S\ref{sec:54-research03-methodology} proves that this selection procedure minimises an upper bound on the expected generalization loss of a model trained on $\tilde{\mathcal{D}}$.
\end{foldable}
